\section{Introduction}
\label{sec:intro}

This paper asks whether any nearby star is transmitting a narrow radio
signal that ALMA happened to record while observing something else. The
public archive holds thousands of hours of millimetre and submillimetre
observations taken for quite different reasons, mostly studies of dust
discs and planet formation, and every one of them also contains a spectrum
of the star at the centre of the field. A transmitter would appear in such
a spectrum as a spike confined to one or two channels. It would sit at the
star's position and nowhere else in the field, and reappear if the star
were observed again. We searched the archive for exactly that. Nothing of
the kind survived.

A technological civilisation may imprint the electromagnetic spectrum
detectably whether or not it intends to, a \emph{technosignature}
\citep{Sheikh2020,SETIreview2025}, and millimetre wavelengths are an
attractive place to look for one. Interstellar scattering and dispersion
fall steeply with frequency, so a high-frequency carrier arrives less
distorted; a given aperture delivers higher gain, and so more radiated
power per watt supplied to the transmitter; and coherent millimetre
emission is already engineered on Earth, for directional links, radar and
beamed power. The band is correspondingly hard to search. The atmosphere
transmits unevenly and adds phase noise, receivers are noisier than at
centimetre wavelengths, and archived correlator channels are thousands of
times coarser than a dedicated narrowband backend, so a carrier is diluted
within one channel and its drift can smear it across several. These bands
are also crowded with molecular lines, a dense background of real
astrophysical signals that a threshold search must separate from an
artificial one.

Six decades of radio searches have examined many thousands of stars, yet
sample only \HaystackFrac{} of the space of transmitter frequency,
bandwidth, power and duty cycle \citep{Wright2018}. Essentially all of that
effort lies below ${\sim}30$\,GHz. Above it we are aware of one published
search, in ALMA Band~3 \citep{Mason2025}, whose targets lie at
\MasonNearKpc\,kpc and beyond; no search above ${\sim}30$\,GHz has yet
reached a star within 40\,pc at all. Few searches at any frequency
reprocess existing interferometric archives, although those archives
already hold a spectrum of every target they have ever pointed at
(\S\ref{sec:background}). That is the regime this paper enters.

The signal sought here is an \emph{unresolved spectral carrier}. It is
emission confined to a few channels, unresolved at the stellar position,
and persistent enough to recur between epochs. Each qualifier is a real
restriction. A carrier broader than the channel width would not be
recovered, nor one that fills the synthesised beam, nor one that appears
once and never again. Its frequency is allowed to drift, as it must for a
transmitter carried on a rotating, orbiting body, and the drift searched is
bounded by a line-of-sight acceleration rather than by a frequency
(\S\ref{sec:statistic}). The fine channels used are
\ChanLoA--\ChanHiA\,kHz wide, so ``narrowband'' here describes the assumed
transmitter and never the resolving power of the data.

One limit on interpretation holds throughout. The targets were chosen by
archival disc and planet-formation programmes, not for this search, and
the sample inherits their preferences in distance and in spectral type
(\S\ref{sec:sample}). Every statistical statement here is therefore about
the ALMA archive-selected stellar sample and not about the solar
neighbourhood.
